\begin{frame}{From Sparse Codes to Joint Integer Tokens}
\small

\textbf{One token identifies an atom and its coefficient.}
For $\boldsymbol d_1=(1,0)^\top$ and
$\boldsymbol d_2=(0,1)^\top$, use the same lookup table
at every residual step:

\medskip
\begin{columns}[T,onlytextwidth]
\begin{column}{0.49\textwidth}
\centering
\renewcommand{\arraystretch}{1.2}
\begin{tabular}{ccl}
\hline
Token & Scaled atom & Stored vector \\
\hline
0 & $\boldsymbol 0$        & $(0,0)^\top$ \\
1 & $-\boldsymbol d_1$     & $(-1,0)^\top$ \\
\textbf{2} & $0.5\boldsymbol d_1$ & $(0.5,0)^\top$ \\
3 & $\boldsymbol d_1$      & $(1,0)^\top$ \\
\textbf{4} & $-\boldsymbol d_2$   & $(0,-1)^\top$ \\
5 & $0.5\boldsymbol d_2$   & $(0,0.5)^\top$ \\
6 & $\boldsymbol d_2$      & $(0,1)^\top$ \\
\hline
\end{tabular}
\end{column}

\begin{column}{0.48\textwidth}
\textbf{Encoding: choose, subtract, repeat.}

For $\boldsymbol z_e=(0.5,-1)^\top$:
\begin{align*}
\boldsymbol r_0 &= (0.5,-1)^\top,\\
\text{choose token }4:\quad
\boldsymbol r_1
&= \boldsymbol r_0-(0,-1)^\top\\
&= (0.5,0)^\top,\\
\text{choose token }2:\quad
\boldsymbol r_2
&= \boldsymbol r_1-(0.5,0)^\top
= \boldsymbol 0.
\end{align*}
The token sequence is $\boxed{[4,2]}$.

\medskip
\textbf{Decoding: look up and add.}
\[
\boldsymbol z_{dl}
= (0,-1)^\top+(0.5,0)^\top
= (0.5,-1)^\top.
\]
\end{column}
\end{columns}

\medskip
\textbf{The transformer predicts one integer at each residual step.}\\
This toy example uses two steps; our model uses
\textbf{four steps} and a shared vocabulary of \textbf{49,153 IDs}.

\end{frame}
